% DEFINED:   sec:interp-closing
% RESOLVED:  sec:three-axes
% FORWARD:   Chapter 8 (reinforcement learning; prose-only, no \Cref)

\section{Closing: the inductive-bias frame, completed}\label{sec:interp-closing}

The book has organized everything around three axes (\Cref{sec:three-axes}):

\begin{enumerate}
\item \textbf{Architecture}: structural priors about how features compose, such
as locality, recurrence, and content-addressable lookup.
\item \textbf{Output head}: priors about the distribution of the response, such
as Bernoulli, Categorical, Gaussian, and Poisson.
\item \textbf{Implementation realization}: whether the gradient procedure can
actually discover the right parameters, given initialization, optimizer, data,
and schedule.
\end{enumerate}

This chapter has been the \textbf{interpretability lens} that sits across all
three. Once a model is trained, mechanistic interpretation tells us \emph{which}
of the bias-permitted solutions the gradient procedure actually picked, and how
legibly. At $M = 8$ the answer was a single, clean, fully readable circuit:
layer 1 aggregates the address, layer 2 dereferences, and ablations confirm both
are load-bearing. At $M = 32$, one increment of scale later, the answer was a
causally perfect pointer whose mechanism the attention maps render only partially
legible, and which a causal trace makes clean again.

That progression is the real shape of the field: legible circuits in the smallest
models, partial legibility as scale grows, and causal probes rather than
weight-reading as the reliable ground truth. The meta-lesson is the one the
$M = 32$ model forced on us. Attention weight is suggestive, not probative; it is
not information flow; and a computation spread across heads, layers, and the
value pathway will hide from it. Reading a model is not staring at its attention
maps. It is intervening on its activations and watching what changes.
Interpretability is how we discover which solution gradient descent picked, and
the $M = 8$ to $M = 32$ step is a warning about how quickly "which" gets hard to
answer.

The closing chapter changes the setting entirely. Until now the learning signal
has been a per-example label, a target the network is trained to match. The next
chapter replaces it with a scalar reward over whole trajectories and turns to
reinforcement learning, where the supervised toolkit assembled across the first
seven chapters reappears as the inner loop of something larger. The
inductive-bias frame carries over, the interpretability lens carries over, and
the mechanism we just dissected becomes one tool among several for understanding
what a learned agent has actually learned.
